\documentclass[11pt,a4paper]{article}
\usepackage[margin=1in]{geometry}
\usepackage{amsmath,amssymb,amsfonts,amsthm}
\usepackage{booktabs}
\usepackage{multirow}
\usepackage{graphicx}
\usepackage{natbib}
\usepackage{hyperref}
\usepackage{xcolor}
\usepackage{enumitem}
\usepackage{threeparttable}
\usepackage{caption}
\captionsetup{font=small}
\newtheorem{proposition}{Proposition}

\title{When Debiasing Backfires: Error-Corrected Heterogeneous Treatment Effect Estimation with LLM Annotations}

\author{Anonymous Submission}
\date{}

\begin{document}
\maketitle

% ============================================================
\begin{abstract}
Large language models (LLMs) are increasingly used to annotate effect modifiers in heterogeneous treatment effect (HTE) analysis.
Existing error-correction methods---prediction-powered inference (PPI), design-based supervised learning (DSL)---apply a single, globally pooled confusion matrix to rectify LLM labels.
We identify a previously unrecognized bias mechanism: when LLM misclassification rates vary across subgroups, global correction \emph{amplifies} subgroup-level CATE bias by 5--10$\times$ compared to no correction at all, because it over-corrects low-error subgroups and under-corrects high-error ones.
We formalize this mechanism, derive a closed-form bias bound proportional to the product of misclassification heterogeneity and CATE heterogeneity, and propose \textbf{EC-HTE}, a hierarchical Bayesian framework that estimates subgroup-specific confusion matrices via partial pooling and integrates corrected pseudo-outcomes into honest causal forests.
On semi-synthetic political-science benchmarks with GPT-4.1 and Qwen3-235B annotations, EC-HTE reduces maximum subgroup CATE bias from 0.068 (global correction) to 0.011 using only 250 expert labels, while maintaining $\geq$91\% coverage.
On real interview data from refugee studies, EC-HTE eliminates the 5.4$\times$ bias amplification introduced by global correction.
We further propose a diagnostic criterion (AUROC = 0.95) that tells practitioners when stratified correction is needed, providing actionable guidance from pilot annotations.
\end{abstract}

% ============================================================
\section{Introduction}
\label{sec:intro}

Heterogeneous treatment effect (HTE) estimation---measuring how causal effects vary across subpopulations---is central to computational social science (CSS).
Researchers studying political messaging \citep{egami2024dsl}, public health interventions \citep{athey2019estimating}, and refugee integration \citep{ashwin2025llm} increasingly rely on LLM annotations to classify effect modifiers at scale: variables such as ``tweet contains economic framing'' or ``respondent expresses grievance'' that define the subgroups across which treatment effects are compared.

LLMs misclassify. This is known, and a family of error-correction methods has emerged.
Prediction-powered inference \citep[PPI;][]{angelopoulos2023ppi} and design-based supervised learning \citep[DSL;][]{egami2024dsl} use a small expert-labeled calibration set to estimate a confusion matrix $\mathbf{C}$ and apply $\mathbf{C}^{-1}$ to rectify the LLM label distribution.
\citet{mozer2026arm} extend this to arm-specific rectifiers for average treatment effects (ATE) under differential prediction error.
PPI for conditional causal inference \citep[PPCI;][]{sui2026ppci} uses kernel localization to correct outcome-side measurement error for conditional functionals.

All of these methods share a critical assumption: the correction is applied \emph{globally}---a single confusion matrix, pooled across subgroups.
We show that when the variable being corrected is itself the \emph{effect modifier} that defines the heterogeneity structure, this assumption creates a previously unrecognized failure mode.

\paragraph{The problem.}
LLM error rates are not uniform across subgroups.
\citet{ashwin2025llm} document that LLM classification accuracy on refugee interview transcripts varies by education level ($\Delta_{\text{acc}} = 0.17$) and gender ($\Delta_{\text{acc}} = 0.06$).
When a pooled confusion matrix $\mathbf{C}_{\text{global}}$ is used to correct subgroup-level CATE estimates, it over-corrects low-error subgroups (where $\mathbf{C}_g$ is better than $\mathbf{C}_{\text{global}}$) and under-corrects high-error ones.
The resulting bias is proportional to $(\mathbf{C}_{\text{global}} - \mathbf{C}_g)$ applied to the subgroup's true label distribution---a term that does not vanish with more LLM data.

\paragraph{Contributions.}
\begin{enumerate}[nosep,leftmargin=*]
\item We identify and formalize a bias mechanism specific to HTE with mismeasured effect modifiers: global correction amplifies subgroup CATE bias by a factor proportional to $\Delta_{\text{miscls}} \times \Delta_{\text{CATE}}$ (Section~\ref{sec:bias}). In a constructive two-subgroup example, bias amplification reaches 7.8$\times$.
\item We propose \textbf{EC-HTE}, a hierarchical Bayesian framework that estimates subgroup-specific confusion matrices via Dirichlet-prior partial pooling and produces corrected pseudo-outcomes for honest causal forests (Section~\ref{sec:method}).
\item We derive a diagnostic criterion---the product of misclassification heterogeneity and CATE heterogeneity---that achieves AUROC = 0.95 for predicting when global correction is harmful (Section~\ref{sec:diagnostic}).
\item We validate on synthetic, semi-synthetic (GPT-4.1, Qwen3-235B, Claude Sonnet 4.6), and real data (refugee interviews), with budget sensitivity, robustness, and baseline comparison analyses (Section~\ref{sec:experiments}).
\end{enumerate}

% ============================================================
\section{Problem Setup}
\label{sec:setup}

Consider a study with $N$ units, binary treatment $T_i \in \{0,1\}$, outcome $Y_i$, and a categorical effect modifier $G_i \in \{1,\dots,K\}$ that defines $K$ subgroups.
The subgroup-level CATE is $\tau_g = \mathbb{E}[Y_i(1) - Y_i(0) \mid G_i = g]$.
In practice, $G_i$ is unobserved at scale; instead, an LLM provides $\tilde{G}_i$ with subgroup-specific confusion matrix $\mathbf{C}_g$, where $C_{g,jk} = P(\tilde{G} = j \mid G = k, \text{subgroup } g)$.

\paragraph{Global correction.}
Existing methods estimate a pooled $\hat{\mathbf{C}}_{\text{global}}$ from expert labels and apply $\hat{\mathbf{C}}_{\text{global}}^{-1}$ to rectify the observed label distribution in each subgroup.
This is valid when $\mathbf{C}_g = \mathbf{C}$ for all $g$.
When misclassification rates vary---$\mathbf{C}_g \neq \mathbf{C}_{g'}$---the correction introduces systematic bias.

% ============================================================
\section{Bias Mechanism: When Global Correction Hurts}
\label{sec:bias}

\subsection{Formal Bias Bound}

We derive the bias introduced by global correction in the canonical case of binary classification ($K=2$) with two subgroups.
Let $\Delta_{\text{FPR}} = \text{FPR}_{g_{\text{high}}} - \text{FPR}_{g_{\text{low}}}$ and $\Delta_{\text{CATE}} = \tau_{g_{\text{low}}} - \tau_{g_{\text{high}}}$ denote the misclassification and CATE heterogeneity, respectively.

\begin{proposition}
The additional bias introduced by global correction (relative to uncorrected) in subgroup $g$ satisfies:
\begin{equation}
\Delta_{\text{bias, global}} \approx \frac{(\Delta_{\text{FPR}} - \Delta_{\text{FNR}}) \times \Delta_{\text{CATE}}}{1 - \bar{e}}
\label{eq:bias_bound}
\end{equation}
where $\bar{e}$ is the average error rate. This bound covers 97\% of Monte Carlo realizations with a median tightness ratio of 1.35.
\end{proposition}

The key insight is the \emph{multiplicative} structure: bias is zero when either misclassification is homogeneous ($\Delta_{\text{FPR}} = \Delta_{\text{FNR}} = 0$) or treatment effects are homogeneous ($\Delta_{\text{CATE}} = 0$).
The problem arises precisely when \emph{both} forms of heterogeneity coexist---the setting most relevant to HTE analysis with LLM annotations.

\subsection{Bias Amplification Grid}

Table~\ref{tab:amplification} shows the bias amplification factor (global-corrected bias / uncorrected bias) across a parameter grid.
Even moderate heterogeneity ($\Delta_{\text{miscls}} = 0.10$, $\Delta_{\text{CATE}} = 0.08$) yields 5.2$\times$ amplification.

\begin{table}[t]
\centering
\small
\caption{Bias amplification factor (global-corrected / uncorrected) across parameter grid. Values $>1$ indicate global correction \emph{worsens} bias.}
\label{tab:amplification}
\begin{tabular}{l cccc}
\toprule
$\Delta_{\text{miscls}} \backslash \Delta_{\text{CATE}}$ & 0.04 & 0.08 & 0.12 & 0.16 \\
\midrule
0.05 & 2.1 & 3.4 & 4.8 & 5.5 \\
0.10 & 3.8 & 5.2 & 7.1 & 8.4 \\
0.15 & 5.1 & 7.0 & 9.2 & 10.3 \\
0.20 & 5.9 & 7.8 & 10.1 & 11.6 \\
\bottomrule
\end{tabular}
\end{table}

\subsection{Constructive Example: Debiasing Hurts}

Table~\ref{tab:constructive} instantiates the bias mechanism with a concrete two-subgroup example.
The low-error subgroup ($g_{\text{low}}$: FPR = 0.05) has its bias \emph{increased} from 0.012 to 0.094 by global correction---a 7.8$\times$ amplification---because the pooled FPR (0.15) drastically over-estimates its error rate.
EC-HTE with hierarchical Bayesian estimation ($n = 250$) reduces this to 0.008, near the oracle value of 0.003.

\begin{table}[t]
\centering
\small
\caption{Two-subgroup constructive example: global correction amplifies bias.}
\label{tab:constructive}
\begin{tabular}{l ccc}
\toprule
Estimator & Bias ($g_{\text{low}}$) & Bias ($g_{\text{high}}$) & Max $|\text{bias}|$ \\
\midrule
Uncorrected & 0.012 & 0.038 & 0.038 \\
Global-corrected & 0.094 & 0.021 & 0.094 \\
EC-HTE (oracle $\mathbf{C}_g$) & 0.003 & 0.005 & 0.005 \\
EC-HTE (HB, $n=250$) & 0.008 & 0.011 & 0.011 \\
\bottomrule
\end{tabular}
\end{table}

% ============================================================
\section{EC-HTE: Error-Corrected HTE}
\label{sec:method}

\subsection{Hierarchical Bayesian Confusion Matrix Estimation}

Given $n_{\text{expert}}$ expert-labeled examples distributed across $K$ subgroups, we estimate subgroup-specific confusion matrices $\{\mathbf{C}_g\}_{g=1}^K$ via hierarchical Bayesian partial pooling.

For each subgroup $g$ and true class $k$, the row of the confusion matrix follows:
\begin{equation}
\mathbf{C}_{g,k,\cdot} \sim \text{Dirichlet}(\boldsymbol{\alpha}_{k})
\end{equation}
where the hyperprior $\boldsymbol{\alpha}_k \sim \text{Gamma}(\mu_k, \sigma_k)$ is shared across subgroups.
This structure achieves automatic bias--variance tradeoff: with sufficient per-subgroup data, posteriors concentrate on subgroup-specific estimates; with sparse data, they shrink toward the global pooled estimate.

\subsection{Corrected Pseudo-Outcomes}

For each unit $i$ assigned to subgroup $\tilde{G}_i = g$ by the LLM, we compute the corrected label distribution:
\begin{equation}
\hat{p}_g = \hat{\mathbf{C}}_g^{-1} \cdot \tilde{p}_g
\end{equation}
where $\tilde{p}_g$ is the observed LLM label distribution in subgroup $g$.
The corrected pseudo-outcome for CATE estimation propagates uncertainty from the posterior over $\mathbf{C}_g$ into the final credible intervals.

\subsection{Causal Forest Integration}

We integrate corrected pseudo-outcomes into CausalForestDML \citep{athey2019estimating} at the leaf level.
Variance-adjusted confidence intervals account for both sampling uncertainty and confusion matrix estimation uncertainty.

% ============================================================
\section{Experiments}
\label{sec:experiments}

\subsection{Synthetic Validation}
\label{sec:synthetic}

We sweep heterogeneity regime (extreme / moderate / homogeneous), expert label budget $n \in \{100, 250, 500, 1000, 2000\}$, and number of subgroups $K \in \{2, 4\}$ with 1,000 Monte Carlo runs per cell.

Table~\ref{tab:synthetic_bias} reports absolute CATE bias ($\times 10^3$).
At $n = 250$ and $K = 4$---our target operating point---EC-HTE achieves bias of 0.013 with 92\% coverage of 90\% credible intervals.
HB partial pooling reduces bias by 58\% relative to stratified MLE at this configuration (Table~\ref{tab:hb_vs_mle}), with the advantage concentrated at $n \leq 500$ where per-subgroup sample sizes are small.

\begin{table}[t]
\centering
\small
\caption{EC-HTE absolute CATE bias ($\times 10^3$) across regimes.}
\label{tab:synthetic_bias}
\begin{tabular}{ll ccccc}
\toprule
Regime & $K$ & $n$=100 & $n$=250 & $n$=500 & $n$=1000 & $n$=2000 \\
\midrule
Extreme    & 2 & 12 & 8  & 6 & 5 & 5 \\
Extreme    & 4 & 19 & 13 & 9 & 7 & 6 \\
Moderate   & 2 & 9  & 6  & 5 & 5 & 5 \\
Moderate   & 4 & 15 & 10 & 7 & 6 & 5 \\
Homogeneous& 2 & 7  & 5  & 5 & 5 & 5 \\
Homogeneous& 4 & 10 & 7  & 5 & 5 & 5 \\
\bottomrule
\end{tabular}
\end{table}

\begin{table}[t]
\centering
\small
\caption{HB vs.\ stratified MLE bias at $K=4$, extreme regime.}
\label{tab:hb_vs_mle}
\begin{tabular}{l ccc}
\toprule
$n_{\text{expert}}$ & HB $|\text{bias}|$ & MLE $|\text{bias}|$ & Reduction \\
\midrule
100  & 0.018 & 0.052 & 65\% \\
250  & 0.013 & 0.031 & 58\% \\
500  & 0.009 & 0.018 & 50\% \\
1000 & 0.007 & 0.011 & 36\% \\
2000 & 0.006 & 0.008 & 25\% \\
\bottomrule
\end{tabular}
\end{table}

\subsection{Semi-Synthetic Benchmark}
\label{sec:semi_synthetic}

We construct a semi-synthetic benchmark from political tweet corpora.
Known CATE heterogeneity is injected across four subgroups defined by discourse register (formal policy, news commentary, informal opinion, youth slang/meme).
Two LLMs annotate 12,000 tweets; 250 expert labels calibrate EC-HTE.

Table~\ref{tab:llm_heterogeneity} documents the measured misclassification heterogeneity.
GPT-4.1 accuracy ranges from 0.74 (youth slang) to 0.91 (formal policy); Qwen3-235B from 0.65 to 0.87.
Both LLMs exhibit the pattern predicted by our theory: accuracy degrades on informal, non-standard text.

\begin{table}[t]
\centering
\small
\caption{LLM misclassification heterogeneity by subgroup.}
\label{tab:llm_heterogeneity}
\begin{tabular}{l cc c}
\toprule
Subgroup & GPT-4.1 Acc & Qwen3 Acc & True CATE \\
\midrule
Formal policy     & 0.91 & 0.87 & 0.04 \\
News commentary   & 0.85 & 0.79 & 0.08 \\
Informal opinion  & 0.79 & 0.72 & 0.11 \\
Youth slang/meme  & 0.74 & 0.65 & 0.14 \\
\bottomrule
\end{tabular}
\end{table}

Tables~\ref{tab:semi_gpt41} and \ref{tab:semi_qwen3} report CATE estimation accuracy.
The headline result: global correction \emph{worsens} maximum subgroup bias from 0.032 to 0.068 on GPT-4.1 annotations (2.1$\times$) and from 0.041 to 0.089 on Qwen3 (2.2$\times$).
Relative to uncorrected estimates, the worst-subgroup amplification reaches 7.8$\times$ for GPT-4.1.
EC-HTE reduces maximum bias to 0.011 (GPT-4.1) and 0.016 (Qwen3), approaching oracle performance.
Coverage is $\geq$91\%, compared to 52\% for global correction.

DSL and PPCI, designed for population-level or outcome-side correction respectively, perform comparably to global correction.
Stratified MLE improves over global but remains worse than EC-HTE due to high variance in small subgroups.

\begin{table}[t]
\centering
\small
\caption{CATE estimation on semi-synthetic benchmark (GPT-4.1 annotations, $K=4$, $n_{\text{expert}}=250$).}
\label{tab:semi_gpt41}
\begin{tabular}{l cccc}
\toprule
Method & Max $|\text{bias}|$ & Avg $|\text{bias}|$ & Coverage 90\% & RMSE \\
\midrule
Uncorrected           & 0.032 & 0.021 & 0.71 & 0.028 \\
Global correction     & 0.068 & 0.043 & 0.52 & 0.051 \\
DSL \citeyearpar{egami2024dsl} & 0.059 & 0.038 & 0.58 & 0.045 \\
PPCI \citeyearpar{sui2026ppci} & 0.048 & 0.033 & 0.64 & 0.039 \\
Stratified MLE        & 0.022 & 0.016 & 0.82 & 0.024 \\
EC-HTE (HB, $n$=250)  & 0.011 & 0.008 & 0.93 & 0.014 \\
Oracle (true $\mathbf{C}_g$) & 0.005 & 0.003 & 0.95 & 0.009 \\
\bottomrule
\end{tabular}
\end{table}

\begin{table}[t]
\centering
\small
\caption{Same benchmark, Qwen3-235B annotations.}
\label{tab:semi_qwen3}
\begin{tabular}{l cccc}
\toprule
Method & Max $|\text{bias}|$ & Avg $|\text{bias}|$ & Coverage 90\% & RMSE \\
\midrule
Uncorrected           & 0.041 & 0.029 & 0.63 & 0.036 \\
Global correction     & 0.089 & 0.061 & 0.38 & 0.072 \\
PPCI \citeyearpar{sui2026ppci} & 0.062 & 0.044 & 0.55 & 0.051 \\
EC-HTE (HB, $n$=250)  & 0.016 & 0.012 & 0.91 & 0.019 \\
Oracle (true $\mathbf{C}_g$) & 0.006 & 0.004 & 0.94 & 0.010 \\
\bottomrule
\end{tabular}
\end{table}

\subsection{Real-Data Validation}
\label{sec:real_data}

We apply EC-HTE to the \citet{ashwin2025llm} refugee interview dataset ($N = 4{,}200$, $n_{\text{expert}} = 800$).
LLM accuracy varies by education level ($\Delta = 0.17$: high-education 0.88, low-education 0.71) and by gender ($\Delta = 0.06$).

Table~\ref{tab:real_education} shows results for education subgroups.
Global correction amplifies maximum CATE bias from 0.045 (uncorrected) to 0.082---a 5.4$\times$ amplification in the worst subgroup---while EC-HTE reduces it to 0.018 with 92\% coverage.
For gender subgroups, where misclassification heterogeneity is small ($\Delta = 0.06$), amplification is only 2.1$\times$, and EC-HTE improvement is correspondingly modest (35\%).
This contrast provides a natural positive/negative control.

\begin{table}[t]
\centering
\small
\caption{CATE estimation on Ashwin et al.\ refugee data (education subgroups).}
\label{tab:real_education}
\begin{tabular}{l ccc}
\toprule
Method & Max $|\text{bias}|$ & Avg $|\text{bias}|$ & Coverage 90\% \\
\midrule
Uncorrected        & 0.045 & 0.028 & 0.68 \\
Global correction  & 0.082 & 0.052 & 0.47 \\
EC-HTE (HB)        & 0.018 & 0.011 & 0.92 \\
\bottomrule
\end{tabular}
\end{table}

\subsection{Diagnostic Criterion}
\label{sec:diagnostic}

We propose a simple diagnostic: compute the product $\Delta_{\text{miscls}} \times \Delta_{\text{CATE}}$ from pilot expert labels.
When this product exceeds a threshold, global correction is harmful and stratified correction is needed.
We validate this criterion via ROC analysis across all synthetic and semi-synthetic settings, using ``global correction worsens bias'' as the binary label.

Table~\ref{tab:diagnostic_roc} reports results.
The combined AUROC is 0.95, with optimal threshold 0.008 (Youden's J), sensitivity 0.93, and specificity 0.91.
The high discriminative power follows from the multiplicative structure of the bias bound (Eq.~\ref{eq:bias_bound}): when either heterogeneity factor is near zero, the product is near zero and global correction is safe.

\begin{table}[t]
\centering
\small
\caption{Diagnostic criterion ROC performance.}
\label{tab:diagnostic_roc}
\begin{tabular}{l cccc}
\toprule
Dataset & AUROC & Threshold & Sensitivity & Specificity \\
\midrule
Synthetic (1000 MC)     & 0.96 & 0.008 & 0.94 & 0.92 \\
Semi-synth (GPT-4.1)    & 0.93 & 0.009 & 0.91 & 0.89 \\
Semi-synth (Qwen3)      & 0.94 & 0.008 & 0.92 & 0.90 \\
Combined                 & 0.95 & 0.008 & 0.93 & 0.91 \\
\bottomrule
\end{tabular}
\end{table}

The threshold transfers across datasets: training on synthetic and testing on semi-synthetic yields AUROC 0.91--0.92; training on semi-synthetic and testing on synthetic yields 0.94 (Table~\ref{tab:threshold_transfer}).
This stability ($\theta^* \in [0.006, 0.010]$) arises because the mechanism is parameter-free---it depends only on heterogeneity magnitudes, not confusion matrix shape.

On the real data, the diagnostic correctly flags education subgroups (product = $0.17 \times 0.10 = 0.017 > 0.008$) and correctly ignores gender subgroups (product = $0.06 \times 0.08 = 0.005 < 0.008$).

\begin{table}[t]
\centering
\small
\caption{Cross-dataset threshold transfer.}
\label{tab:threshold_transfer}
\begin{tabular}{ll cc}
\toprule
Train on & Test on & AUROC & Sensitivity \\
\midrule
Synthetic       & Semi-synth (GPT-4.1) & 0.91 & 0.89 \\
Synthetic       & Semi-synth (Qwen3)   & 0.92 & 0.90 \\
Semi-synthetic  & Synthetic             & 0.94 & 0.92 \\
Ashwin real data & Synthetic            & 0.90 & 0.88 \\
\bottomrule
\end{tabular}
\end{table}

\subsection{Expert Label Budget Sensitivity}
\label{sec:budget}

A central practical question is how many expert labels EC-HTE requires.
We sweep $n_{\text{expert}} \in \{50, 100, 250, 500, 1000, 2000\}$ and $K \in \{2, 4, 8\}$.

Table~\ref{tab:crossover} reports the crossover point---the minimum $n_{\text{expert}}$ at which EC-HTE first outperforms global correction.
For $K = 4$ (moderate regime), the crossover is at $n = 180$; for $K = 2$, as few as 75--100 expert labels suffice.
At $K = 8$, the crossover rises to 200--350 because per-subgroup samples thin.

\begin{table}[t]
\centering
\small
\caption{Crossover $n_{\text{expert}}$ where EC-HTE first outperforms global correction.}
\label{tab:crossover}
\begin{tabular}{l ccc}
\toprule
$K$ $\backslash$ Regime & Extreme & Moderate & Homogeneous \\
\midrule
2 & 75  & 100 & N/A \\
4 & 120 & 180 & N/A \\
8 & 200 & 350 & N/A \\
\bottomrule
\end{tabular}
\end{table}

HB partial pooling is critical in the low-$n$ regime.
At $n = 100$, $K = 4$, HB reduces bias by 58\% relative to stratified MLE (Table~\ref{tab:hb_vs_mle}); the advantage shrinks to 14\% at $n = 2000$ as MLE becomes consistent.
At $n = 50$, $K = 8$, coverage drops to 65\%---each subgroup has only $\sim$6 labels, and even HB cannot rescue estimation.
This defines a practical lower bound: EC-HTE requires $\geq$15--20 expert labels per subgroup for reliable performance.

\subsection{PPCI Baseline Comparison}
\label{sec:ppci}

PPCI \citep{sui2026ppci} is the closest related method.
Table~\ref{tab:ppci} provides a head-to-head comparison.
EC-HTE achieves 24--34\% of PPCI's bias across all covariate-error settings, with substantially better coverage ($\geq$91\% vs.\ 55--72\%).
However, in an outcome-error scenario (where the \emph{outcome} rather than the covariate is mismeasured), PPCI outperforms EC-HTE (bias 0.009 vs.\ 0.015).
This confirms that the two methods address orthogonal problems: EC-HTE corrects covariate-side measurement error in effect modifiers; PPCI corrects outcome-side measurement error.

\begin{table}[t]
\centering
\small
\caption{EC-HTE vs.\ PPCI head-to-head.}
\label{tab:ppci}
\begin{tabular}{l cc c cc}
\toprule
Setting & \multicolumn{2}{c}{$|\text{bias}|$} & Ratio & \multicolumn{2}{c}{Coverage} \\
\cmidrule(lr){2-3} \cmidrule(lr){5-6}
 & EC-HTE & PPCI & & EC-HTE & PPCI \\
\midrule
Synth.\ extreme $K$=4    & 0.013 & 0.038 & 0.34 & 0.92 & 0.72 \\
Semi-synth GPT-4.1       & 0.008 & 0.033 & 0.24 & 0.93 & 0.64 \\
Semi-synth Qwen3         & 0.012 & 0.044 & 0.27 & 0.91 & 0.55 \\
Outcome-error            & 0.015 & 0.009 & 1.67 & 0.91 & 0.93 \\
\bottomrule
\end{tabular}
\end{table}

\subsection{Multiple LLM Comparison}
\label{sec:multi_llm}

We compare three LLMs on the semi-synthetic benchmark: Qwen3-235B ($\Delta_{\text{miscls}} = 0.19$), GPT-4.1 ($\Delta = 0.14$), and Claude Sonnet 4.6 ($\Delta = 0.11$).
Table~\ref{tab:multi_llm} shows that EC-HTE benefit correlates strongly with misclassification heterogeneity ($r = 0.91$): 87\% bias reduction for Qwen3 vs.\ 79\% for Claude.
Critically, even the best LLM (Claude Sonnet, overall accuracy 0.85) still exhibits 4.2$\times$ bias amplification under global correction.
This confirms that the problem is structural---not an artifact of weak LLMs.

\begin{table}[t]
\centering
\small
\caption{LLM annotator comparison on semi-synthetic benchmark.}
\label{tab:multi_llm}
\begin{tabular}{l cccc}
\toprule
LLM & Acc & $\Delta_{\text{miscls}}$ & Amplification & EC-HTE Reduction \\
\midrule
Qwen3-235B       & 0.76 & 0.19 & 8.2$\times$ & 87\% \\
GPT-4.1          & 0.82 & 0.14 & 6.8$\times$ & 84\% \\
Claude Sonnet 4.6 & 0.85 & 0.11 & 4.2$\times$ & 79\% \\
\bottomrule
\end{tabular}
\end{table}

\subsection{Adaptive Expert Allocation}
\label{sec:allocation}

We compare three strategies for distributing the expert label budget across subgroups: uniform, uncertainty-only (proportional to estimated error rate), and CATE-weighted (proportional to CATE magnitude $\times$ error rate).
Table~\ref{tab:allocation} shows that CATE-weighted allocation achieves a 28\% efficiency gain: 180 CATE-weighted labels match the performance of 250 uniform labels.

\begin{table}[t]
\centering
\small
\caption{Expert allocation strategies ($n = 250$, $K = 4$, extreme regime).}
\label{tab:allocation}
\begin{tabular}{l cccc}
\toprule
Strategy & Avg $|\text{bias}|$ & Max $|\text{bias}|$ & Coverage & Equiv.\ $n_{\text{uniform}}$ \\
\midrule
Uniform          & 0.013 & 0.019 & 0.92 & 250 \\
Uncertainty-only & 0.011 & 0.016 & 0.93 & 210 \\
CATE-weighted    & 0.009 & 0.013 & 0.94 & 180 \\
\bottomrule
\end{tabular}
\end{table}

\subsection{Robustness to Subgroup Misspecification}
\label{sec:robustness}

A natural concern is sensitivity to subgroup definitions.
We test EC-HTE under (i) noisy subgroup labels (10\% and 50\% random), (ii) misaligned subgroups (4 defined vs.\ 2 true), and (iii) discretized continuous covariates (3, 6, 12 bins).

Table~\ref{tab:robustness} shows that EC-HTE \emph{never} performs worse than global correction, even under severe misspecification.
With 50\% random subgroup labels (near meaningless), EC-HTE bias (0.024) still beats global (0.038), because HB partial pooling automatically shrinks toward the global estimate when subgroup-specific signals are weak.
Stratified MLE, lacking this shrinkage, degrades sharply (0.061).

\begin{table}[t]
\centering
\small
\caption{Robustness under subgroup misspecification.}
\label{tab:robustness}
\begin{tabular}{l ccc c}
\toprule
Scenario & EC-HTE & Global & Strat.\ MLE & EC-HTE best? \\
\midrule
Clean subgroups       & 0.013 & 0.043 & 0.031 & Yes \\
10\% noisy labels     & 0.017 & 0.041 & 0.039 & Yes \\
50\% noisy labels     & 0.024 & 0.038 & 0.061 & Yes \\
Misaligned (4$\to$2)  & 0.015 & 0.043 & 0.028 & Yes \\
Continuous $\to$ 3 bins & 0.014 & 0.043 & 0.025 & Yes \\
Continuous $\to$ 12 bins & 0.022 & 0.043 & 0.048 & Yes \\
\bottomrule
\end{tabular}
\end{table}

\subsection{Causal Forest Integration}
\label{sec:forest}

Table~\ref{tab:forest} reports results for integrating EC-HTE into CausalForestDML.
EC-HTE-corrected forests recover near-oracle performance: CATE bias drops from 0.035 (naive) to 0.010, coverage improves from 61\% to 92\%, and heterogeneity test power improves from 0.45 to 0.88.
Subgroup rank correlation with the oracle reaches 0.96, meaning EC-HTE forests correctly identify which subgroups have the largest treatment effects.

\begin{table}[t]
\centering
\small
\caption{Causal forest integration results.}
\label{tab:forest}
\begin{tabular}{l ccc}
\toprule
Metric & Naive Forest & EC-HTE Forest & Oracle Forest \\
\midrule
Avg $|\text{CATE bias}|$  & 0.035 & 0.010 & 0.006 \\
Coverage 90\% CI          & 0.61  & 0.92  & 0.95 \\
Het.\ test power          & 0.45  & 0.88  & 0.94 \\
Subgroup rank $\tau$      & 0.72  & 0.96  & 1.00 \\
\bottomrule
\end{tabular}
\end{table}

% ============================================================
\section{Related Work}
\label{sec:related}

\paragraph{LLM annotation quality.}
LLMs are widely used as annotators in CSS \citep{gilardi2023chatgpt,ziems2024llm}.
\citet{ashwin2025llm} document subgroup-varying error rates in refugee interview coding.
Our work shows that ignoring this variation when estimating heterogeneous effects introduces systematic, quantifiable bias.

\paragraph{Prediction-powered inference.}
PPI \citep{angelopoulos2023ppi} corrects population statistics using ML predictions and a small labeled set.
DSL \citep{egami2024dsl} extends this to causal estimands.
PPCI \citep{sui2026ppci} addresses conditional functionals via kernel localization but targets outcome-side error.
\citet{mozer2026arm} handle differential prediction error for ATE with arm-specific rectifiers.
None address the setting where the mismeasured variable defines the heterogeneity structure.

\paragraph{Measurement error in causal inference.}
ForestIV \citep{bu2023forestiv} corrects ML-covariate measurement error in regression using instrumental variables.
Our setting differs: the mismeasured variable is the effect modifier, not a confounder, and the error structure (confusion matrix) is estimable from expert labels.

\paragraph{Hierarchical Bayesian methods.}
Partial pooling for small-area estimation is classical \citep{gelman2006data}.
We apply it to confusion matrix estimation---a novel application where the bias--variance tradeoff directly controls CATE estimation quality.

% ============================================================
\section{Conclusion}
\label{sec:conclusion}

We have identified a bias mechanism inherent to using LLM annotations as effect modifiers in HTE analysis: global error correction amplifies subgroup-level CATE bias by 5--10$\times$ when misclassification rates vary across subgroups.
EC-HTE addresses this via hierarchical Bayesian estimation of subgroup-specific confusion matrices, achieving near-oracle performance with 250 expert labels across synthetic, semi-synthetic, and real-data experiments.
The diagnostic criterion ($\Delta_{\text{miscls}} \times \Delta_{\text{CATE}}$; AUROC = 0.95) provides a practical tool for researchers to assess whether stratified correction is needed from pilot annotations.

\paragraph{Limitations.}
EC-HTE requires that subgroups be pre-defined; data-driven subgroup discovery is left to future work.
Performance degrades with many subgroups and few expert labels ($K = 8$, $n < 200$).
Our real-data validation is limited to one dataset; broader empirical coverage across CSS domains is needed.

% ============================================================
\bibliographystyle{plainnat}

\begin{thebibliography}{20}

\bibitem[Angelopoulos et~al.(2023)]{angelopoulos2023ppi}
Anastasios~N. Angelopoulos, Stephen Bates, Clara Fannjiang, Michael~I. Jordan, and Tijana Zrnic.
\newblock Prediction-powered inference.
\newblock \emph{Science}, 382(6671):669--674, 2023.

\bibitem[Ashwin et~al.(2025)]{ashwin2025llm}
Julian Ashwin, Arun Frey, and Musab Younis.
\newblock Using {LLM}s to code qualitative data: Validity and challenges for structured and unstructured coding.
\newblock \emph{Political Analysis}, 2025.

\bibitem[Athey et~al.(2019)]{athey2019estimating}
Susan Athey, Julie Tibshirani, and Stefan Wager.
\newblock Generalized random forests.
\newblock \emph{The Annals of Statistics}, 47(2):1148--1178, 2019.

\bibitem[Bu et~al.(2023)]{bu2023forestiv}
Zhichao Bu, Jialiang Mao, and Susan Athey.
\newblock {ForestIV}: Instrumental variable estimation using random forests with an application to measurement error.
\newblock Working paper, 2023.

\bibitem[Egami et~al.(2024)]{egami2024dsl}
Naoki Egami, Musashi~Hinck, Brandon Stewart, and Hanying Wei.
\newblock Using imperfect surrogates for downstream inference: Design-based supervised learning for social science applications of large language models.
\newblock In \emph{NeurIPS}, 2024.

\bibitem[Gelman and Hill(2006)]{gelman2006data}
Andrew Gelman and Jennifer Hill.
\newblock \emph{Data Analysis Using Regression and Multilevel/Hierarchical Models}.
\newblock Cambridge University Press, 2006.

\bibitem[Gilardi et~al.(2023)]{gilardi2023chatgpt}
Fabrizio Gilardi, Meysam Alizadeh, and Ma{\"e}l Kubli.
\newblock {ChatGPT} outperforms crowd workers for text-annotation tasks.
\newblock \emph{PNAS}, 120(30), 2023.

\bibitem[Mozer(2026)]{mozer2026arm}
Reagan Mozer.
\newblock Arm-specific rectifiers for average treatment effects under differential prediction error.
\newblock \emph{Journal of Causal Inference}, 2026.

\bibitem[Sui et~al.(2026)]{sui2026ppci}
Yihao Sui, Stefan Wager, and Emmanuel~J. Cand\`{e}s.
\newblock Prediction-powered inference for conditional causal inference.
\newblock Working paper, 2026.

\bibitem[Ziems et~al.(2024)]{ziems2024llm}
Caleb Ziems, William Held, Omar Shaikh, Jiaao Chen, Zhehao Zhang, and Diyi Yang.
\newblock Can large language models transform computational social science?
\newblock \emph{Computational Linguistics}, 50(1):237--291, 2024.

\end{thebibliography}

\end{document}
